\section{Conclusions}
\label{sec:conclusions}

We have constructed a near-field transport model for a single leaking
\gls{WDP} in an access-tunnel segment, explicitly resolving the
bentonite/backfill and EDZ domains. Flow and transport in the surrounding rock
are represented by a mixed-dimensional \gls{DFM} formulation with a
stochastically generated \gls{DFN}. A complementary hydro-mechanical analysis
was used to estimate permeability changes induced  by excavation and
\gls{DGR} closure. It was used to define uncertainty ranges for EDZ parameters.
These model components were combined with a stochastic description of key
inputs, enabling a variance-based sensitivity analysis using Sobol indices for
parameter groups.

The transport model is diffusion-dominated for the considered parameter ranges.
Locally, advection can become more pronounced along major fractures and along
short ``matrix bridges'' connecting fractures, but sustained
advection-dominated behaviour is unlikely even under substantially higher
velocities because increasing velocities also increase hydrodynamic dispersion.

The DFN has a decisive effect on breakthrough. In the nominal fractured
scenario (Case~0), the characteristic breakthrough time is on the order of
$2\,\munit{ky}$, but with very large variability, ranging from about
$0.11\,\munit{ky}$ to several $100\,\munit{ky}$. In contrast, the DFN-free
control case (Case~1) delays breakthrough to roughly $100\,\munit{ky}$ and
exhibits much smaller variability. Consistent with progressive mixing and the
diminishing influence of early-time pathways, the variance of the time-dependent
\gls{QoI} $I(t)$ decreases with time; we therefore also interpret results using
the aggregated scalar quantity $\max I(t)$, which captures the peak
concentration over the full simulation horizon.

The sensitivity analysis shows that, once fractures are present, the dominant
contributions to the output variance come from the realized fracture geometry
and the DFN parameterization. In Case~0, the DFN realization and DFN parameter
group dominate the total-effect indices, while first-order effects remain small
and the main influence arises through interactions. The strongest and most
robust interaction is between the DFN realization and the DFN parameter group,
implying that variance reduction requires improving DFN characterization;
however, simply fitting DFN parameters to larger outcrop datasets may be
counterproductive unless the nature of the uncertainty is clarified
(epistemic vs.\ spatial variability vs.\ aleatory). Numerical uncertainty due
to meshing is smaller, but not entirely negligible.

The DFN-free control analysis in Case~1 shows a qualitatively different
sensitivity pattern: without fractures, interaction effects and meshing
uncertainty are weak, and the dominant contributions come from the bulk-rock
and EDZ parameter groups, while the backfill contribution is the least
significant. Separating DFN uncertainty into geometric descriptors
(\emph{DFN population}) and transmissivity-related parameters
(\emph{DFN transport}) in Case~2 shows that geometry remains more influential
than fracture transmissivity, but also reveals a small visible contribution of
the DFN transport group, on the order of $5~\%$ of the variance. Since the
uncertainty assigned to fracture transmissivity is based only on a rough
estimate, its role may still be underestimated.

Case~3 explores the possibility that the background rock conductivity in
Case~0 is underestimated because unresolved small-scale fractures are not
represented explicitly. Under elevated bulk-rock conductivity and dispersivity,
the concentration variability increases and the dominant control shifts further
toward surrounding-rock connectivity and realized fracture geometry. In this
regime, the backfill and EDZ parameters mainly perturb already efficient
pathways, whereas among the material groups only the bulk-rock parameters
remain clearly significant.

Finally, while the same workflow is transferable to regional-scale transport
analyses, the dominant sensitivities may change substantially, in particular
due to the expected much larger uncertainty and heterogeneity of rock-matrix
properties at regional scales.
